% (C) 2026 Ralf Stephan, in collaboration with Claude Code.
% Released under CC0 1.0 Universal (public-domain dedication).
% See https://creativecommons.org/publicdomain/zero/1.0/

\documentclass[11pt]{article}
\usepackage[T1]{fontenc}
\usepackage[utf8]{inputenc}
\usepackage{lmodern}
\usepackage[margin=2.5cm]{geometry}
\usepackage{amsmath,amssymb}
\usepackage{booktabs,tabularx}
\usepackage{braket}
\usepackage[hidelinks]{hyperref}

% Lean names and paths: verbatim, breakable at . _ /
\DeclareUrlCommand\lean{\urlstyle{tt}}
\def\UrlBreaks{\do\.\do\_\do\/}
\newcolumntype{L}{>{\raggedright\arraybackslash}X}
\newcolumntype{B}{>{\raggedright\arraybackslash}p{0.42\textwidth}}
\emergencystretch=6em

\title{Lean formalization of an Assumptions of Physics technical brief:\\ log}
\author{}
\date{}

\begin{document}
\maketitle

\noindent
This log records which statements of the technical brief in this directory were checked in Lean,
where the formal version differs from the brief, and any problems found in the brief.
The brief is available at \url{https://assumptionsofphysics.org/resources/briefs}.

\section*{Setup}

\begin{itemize}
\item Lean \texttt{v4.36.0-rc2}; Mathlib at commit \texttt{1b1ccf4b}; physlib (\texttt{master})
  and TauCeti (commit \texttt{e7e7f52}, on the same Mathlib) as further dependencies.
\item Build: \texttt{lake build} in the project root.
\item Every file is checked with \texttt{\#print axioms} on its main results. A result is
  recorded here only if it depends on no axioms beyond \lean{propext},
  \lean{Classical.choice} and \lean{Quot.sound}, and so has no \lean{sorry}.
\item Imports are kept to the Mathlib modules actually needed.
\item Material that belongs in Mathlib is kept in \lean{ForMathlib/}, in Mathlib's directory
  layout (library \lean{ForMathlib}), and imported by the brief files.
\end{itemize}


%======================================================================

\section{Inconsistent convergence in Hilbert space}

\paragraph{Files.}\leavevmode
\begin{itemize}
\item File: \lean{briefs/InconsistentConvergence/InconsistentConvergence.lean}
\item Library: \lean{InconsistentConvergence}
\item Imports: \lean{Mathlib.Analysis.Normed.Lp.lpSpace}
\item Second file, for the oscillator: module \lean{InconsistentConvergence.Oscillator}, in the
  subdirectory \lean{InconsistentConvergence/}. It is separate from the main file, so that the
  comparator lane, which certifies the main file, is unchanged. Its imports:
  \begin{itemize}
  \item \lean{InconsistentConvergence};
  \item \lean{ForMathlib.Analysis.Distribution.SchwartzSpace.HermiteReal};
  \item \lean{TauCeti.Analysis.SpecialFunctions.Hermite.Function.Oscillator}.
  \end{itemize}
  The \lean{ForMathlib} file identifies $\mathcal{S}(\mathbb{R})$, with its topology, with the
  rapidly decreasing Hermite series. It transfers the multi-index Hermite expansions on
  $\mathbb{R}^1$, in the same directory, to $\mathbb{R}$.
\item Third file, for the operators: module \lean{InconsistentConvergence.NumberOperator}, in
  the same subdirectory and separate from the main file for the same reason. Its imports:
  \begin{itemize}
  \item \lean{InconsistentConvergence};
  \item \lean{ForMathlib.Analysis.InnerProductSpace.l2Diagonal};
  \item \lean{ForMathlib.LinearAlgebra.LinearPMap}.
  \end{itemize}
  The \lean{ForMathlib} files provide the diagonal operators on $\ell^2$, self-adjoint on their
  maximal domains, and the composition of unbounded operators on its natural domain.
\item Fourth file, for the Schwartz sequence space, in the same subdirectory and separate from
  the main file for the same reason: module \lean{InconsistentConvergence.SchwartzSequences}.
  Its imports:
  \begin{itemize}
  \item \lean{InconsistentConvergence.Oscillator};
  \item \lean{ForMathlib.Analysis.Normed.Lp.SchwartzSeq}.
  \end{itemize}
  The \lean{ForMathlib} file constructs $S(\mathbb{N})$ as a topological vector space.
\item Fifth file, for the observables on $S(\mathbb{N})$, in the same subdirectory and separate
  from the main file for the same reason: module \lean{InconsistentConvergence.SchwartzObservables}.
  Its imports:
  \begin{itemize}
  \item \lean{InconsistentConvergence.NumberOperator};
  \item \lean{InconsistentConvergence.SchwartzSequences}.
  \end{itemize}
  The construction of continuous linear maps of $S(\mathbb{N})$ from weighted bounds, and the
  continuity of their expectations, are in \lean{ForMathlib.Analysis.Normed.Lp.SchwartzSeq}.
\end{itemize}

\paragraph{Setting.} $H = \ell^2(\mathbb{N})$ (\texttt{lp (fun \_ : }$\mathbb{N}$\texttt{ => }$\mathbb{C}$\texttt{) 2}), with
$\ket{n} = \delta_n$ the oscillator eigenstates and
\[
  \ket{e^j} = \sqrt{\tfrac{j-1}{j}}\,\ket{0} + \sqrt{\tfrac{1}{j}}\,\ket{j}.
\]
In the oscillator file, $U$ is the unitary $\ell^2(\mathbb{N}) \to L^2(\mathbb{R})$ of the
Hermite basis (\lean{TauCeti.hermiteHilbertBasis}), with $U\ket{n} = \psi_n$, the $n$-th
Hermite function ($m = \omega = \hbar = 1$). The oscillator file describes $S(\mathbb{N})$ by the
seminorms $\sup_n (n+1)^K |x_n|$, $K \in \mathbb{N}$; the fourth file constructs it as a
topological vector space, with the brief's seminorms.

\paragraph{Results.}
\begin{center}
\begin{tabularx}{\textwidth}{@{}BL@{}}
\toprule
Brief & Lean \\
\midrule
eq.~(2): $\|e^j - \ket{0}\|^2 = \bigl(\sqrt{(j-1)/j} - 1\bigr)^2 + 1/j$ & \lean{norm_sub_ground_sq} \\
$e^j \to \ket{0}$ in $\ell^2$ & \lean{tendsto_e} \\
eq.~(3): $\sum_n n^2 |e^j_n|^2 = j \to \infty$ & \lean{secondMoment_e}, \lean{tendsto_secondMoment_e} \\
the $\ket{j}$ alone have no limit in $\ell^2$ & \lean{not_cauchySeq_level}: $\|\ket{j}-\ket{k}\|^2 = 2$ for $j\ne k$ \\
eq.~(4): $\braket{e^j|N|e^j} = 1$ for $j \ge 1$ & \lean{numberExp_e} \\
eq.~(5), the inconsistency & \lean{not_continuousAt_numberExp}: $\langle N\rangle$ is not continuous at $\ket{0}$ in $\ell^2$ \\
\bottomrule
\end{tabularx}
\end{center}

\paragraph{Results: the oscillator.}
\begin{center}
\begin{tabularx}{\textwidth}{@{}BL@{}}
\toprule
Brief & Lean \\
\midrule
$U\ket{n} = \psi_n$ & \lean{hermiteUnitary_level}; the coordinates of $U^{-1} g$ are
  $\braket{\psi_n|g}$ (\lean{hermiteUnitary_symm_apply}) \\
$\ket{n}$ is the oscillator eigenstate with $n$ quanta & \lean{oscillator_level}:
  $-\psi_n''/2 + x^2\psi_n/2 = (n + \tfrac12)\psi_n$ \\
$S(\mathbb{N})$ is the Schwartz space & \lean{exists_schwartz_iff}: $Ux \in
  \mathcal{S}(\mathbb{R})$ iff $\sup_n (n+1)^K |x_n| < \infty$ for every $K$ \\
and has its topology & \lean{tendsto_schwartz_iff}: $f_j \to g$ in $\mathcal{S}(\mathbb{R})$
  iff $U^{-1}f_j \to U^{-1}g$ in every seminorm of $S(\mathbb{N})$ \\
the $e^j$ are Schwartz states & \lean{exists_schwartz_e}: $Ue^j \in
  \mathcal{S}(\mathbb{R})$ \\
$e^j \to \ket{0}$, but not in $S(\mathbb{N})$ & \lean{tendsto_hermiteUnitary_e}: $Ue^j \to
  \psi_0$ in $L^2(\mathbb{R})$; \lean{not_tendsto_schwartz_e}: $Ue^j$ has no limit in
  $\mathcal{S}(\mathbb{R})$ \\
\bottomrule
\end{tabularx}
\end{center}

\paragraph{Results: the number operator.} In the operator file, $a$ and $a^\dagger$ act by
$(ax)_n = \sqrt{n+1}\,x_{n+1}$, $(a^\dagger y)_0 = 0$ and $(a^\dagger y)_{n+1} = \sqrt{n+1}\,y_n$,
and $N$ by $(Nx)_n = n\,x_n$, each on its maximal domain:
$D(a) = \{x \mid (\sqrt{n}\,x_n)_n \in \ell^2\}$ and $D(N) = \{x \mid (n\,x_n)_n \in \ell^2\}$.
\begin{center}
\begin{tabularx}{\textwidth}{@{}BL@{}}
\toprule
Brief & Lean \\
\midrule
$a\ket{n+1} = \sqrt{n+1}\,\ket{n}$, $a\ket{0} = 0$, $a^\dagger\ket{n} = \sqrt{n+1}\,\ket{n+1}$
  & \lean{annihilation_level_succ}, \lean{annihilation_level_zero}, \lean{creation_level} \\
$a^\dagger$ is the adjoint of $a$ & \lean{adjoint_annihilation} \\
$N$ is self-adjoint, with $N\ket{n} = n\ket{n}$ & \lean{isSelfAdjoint_numberOp},
  \lean{numberOp_level}; domain: \lean{mem_domain_numberOp} \\
$N = a^\dagger a$ & \lean{numberOp_eq_adjoint_compMax}, with the composition on
  $\{x \in D(a) \mid ax \in D(a^\dagger)\}$ \\
$\braket{x|N|x} = \sum_n n\,|x_n|^2$ & \lean{inner_numberOp}, on $D(N)$;
  \lean{numberExp_eq_norm_sq}: $\sum_n n\,|x_n|^2 = \|ax\|^2$ on $D(a)$ \\
eq.~(4) for the operator & \lean{inner_numberOp_e}: $e^j \in D(N)$ and
  $\braket{e^j|N|e^j} = 1$ for $j \ge 1$ \\
\bottomrule
\end{tabularx}
\end{center}

\paragraph{Results: the Schwartz sequence space.} In the fourth file, $S(\mathbb{N})$ is the
type \lean{SchwartzSeq} of sequences with $((n+1)^k x_n)_n \in \ell^2$ for every $k$, with the
topology of the seminorms $\|((n+1)^k x_n)_n\|$.
\begin{center}
\begin{tabularx}{\textwidth}{@{}BL@{}}
\toprule
Brief & Lean \\
\midrule
$S(\mathbb{N}) \subset \ell^2$, with the topology of its seminorms
  & \lean{SchwartzSeq.withSeminorms}; the inclusion into $\ell^2$ is continuous
  (\lean{SchwartzSeq.continuous_toLp}) \\
Prop.~1: $e^j \to a$ in $\ell^2$, but $e^j$ does not converge in $S(\mathbb{N})$
  & \lean{tendsto_toLp_eS}, \lean{not_tendsto_eS}: $\|((n+1)\,e^j_n)_n\| \ge \sqrt{j}$
  (\lean{sqrt_le_seminorm_eS}) \\
each $b^j$ is in $S(\mathbb{N})$, but the sequence has no limit in $S(\mathbb{N})$
  & \lean{levelS}, \lean{not_tendsto_levelS} \\
$S(\mathbb{N})$ is the Schwartz space of the oscillator & \lean{hermiteSchwartzEquiv}:
  $S(\mathbb{N}) \cong \mathcal{S}(\mathbb{R})$ as topological vector spaces, given by $U$
  (\lean{toLp_hermiteSchwartzEquiv}), with $\ket{n} \mapsto \psi_n$
  (\lean{toLp_hermiteSchwartzEquiv_levelS}) \\
\bottomrule
\end{tabularx}
\end{center}

\paragraph{Results: observables on $S(\mathbb{N})$.} In the fifth file, $a$, $a^\dagger$ and $N$
are restricted to $S(\mathbb{N})$, which lies in their domains. The bounds
$(n+1)^k |(ax)_n| \le p_{k+1}(x)$, $(n+1)^k |(a^\dagger x)_n| \le 2^k p_{k+1}(x)$ and
$(n+1)^k |(Nx)_n| \le p_{k+1}(x)$, with $p_k(x) = \|((n+1)^k x_n)_n\|$, make them continuous
linear maps $S(\mathbb{N}) \to S(\mathbb{N})$ (\lean{SchwartzSeq.mkCLM}).
\begin{center}
\begin{tabularx}{\textwidth}{@{}BL@{}}
\toprule
Brief & Lean \\
\midrule
$a$, $a^\dagger$, $N$ act on $S(\mathbb{N})$ & \lean{annihilationS}, \lean{creationS},
  \lean{numberOpS}, agreeing with the operators of the fourth file
  (\lean{toLp_numberOpS}); $N = a^\dagger a$ (\lean{numberOpS_eq_comp}) \\
the Schwartz topology does not allow this kind of convergence: $\langle N\rangle$
  & \lean{continuous_inner_numberOpS}, \lean{continuous_numberExp}: $\braket{x|N|x}$ is
  continuous on $S(\mathbb{N})$ \\
Prop.~1, from the expectation & \lean{not_tendsto_eS_of_numberExp}: a limit of $e^j$ in
  $S(\mathbb{N})$ would be $\ket{0}$, with $\langle N\rangle = 0 \ne 1$ \\
the same for all moments & \lean{continuous_inner_of_mem_adjoin}: $\braket{x|T|x}$ is continuous
  on $S(\mathbb{N})$ for every polynomial $T$ in $a$ and $a^\dagger$, among them
  $\hat{x} = (a + a^\dagger)/\sqrt{2}$, $\hat{p} = i(a^\dagger - a)/\sqrt{2}$ and
  $N + \tfrac12$ (\lean{positionS}, \lean{momentumS}, \lean{hamiltonianS}) \\
\bottomrule
\end{tabularx}
\end{center}

\paragraph{Differences from the brief.}
\begin{itemize}
\item Indexing. Proposition~1 indexes sequences from $1$, with the ground term at $i = 1$. The
  formalization indexes from $0$, matching the oscillator of Section~3. The divergence is the
  same; with the brief's indexing the sum is $(j-1)/j + j$ instead of $j$.
\item No limit in $S(\mathbb{N})$. In the main file, what is proved is that the seminorm
  $\sum_n n^2 |x_n|^2$ is unbounded along $e^j$, which is the brief's argument. That $e^j$ has
  no limit in the topology of $S(\mathbb{N})$ is proved in the fourth file.
\item Seminorms of $S(\mathbb{N})$. The brief's seminorms $\sum_i |i^k x_i|^2$, with $i \ge 1$,
  are the squares of $\|((n+1)^k x_n)_n\|$ with indices from $0$. The oscillator file uses the
  seminorms $\sup_n (n+1)^K |x_n|$, which define the same space and the same topology
  (\lean{SchwartzSeq.mem_submodule_iff}, \lean{SchwartzSeq.tendsto_iff}).
\item In the main file, $\langle N\rangle$ is defined on all of $\ell^2$ as the weighted sum
  $\sum_n n\,|x_n|^2$, which has a default value outside the domain of $N$. The discontinuity
  is still meaningful, because every $e^j$ is finitely supported. The operator file constructs
  $N = a^\dagger a$ and shows that the weighted sum is $\braket{x|N|x}$ on $D(N)$.
\item All moments. The brief says that the Schwartz topology ``keeps track of all moments''.
  The formalization reads this as the continuity on $S(\mathbb{N})$ of the expectations of the
  polynomials in $a$ and $a^\dagger$. The proof does not use the polynomial form: the
  expectation of every continuous linear map $S(\mathbb{N}) \to S(\mathbb{N})$ is continuous
  (\lean{SchwartzSeq.continuous_inner_apply}), and the polynomials are such maps because $a$
  and $a^\dagger$ are.
\end{itemize}

\paragraph{Problems found in the brief.} No mathematical problems. Minor: the distance is
written $d(x, y) = \sum_i (x_i - y_i)^2$, which is the squared distance and, for complex
entries, should read $\sum_i |x_i - y_i|^2$.

\end{document}
